\documentclass[11pt,a4paper]{article}
\usepackage[utf8]{inputenc}
\usepackage[T1]{fontenc}
\usepackage{lmodern,amsmath,amssymb,booktabs,tabularx,array}
\usepackage[margin=2.2cm]{geometry}
\usepackage[hidelinks]{hyperref}
\usepackage{listings}
\lstset{basicstyle=\ttfamily\small,breaklines=true,columns=fullflexible,keepspaces=true}
\setlength{\parindent}{0pt}
\setlength{\parskip}{5pt}
\renewcommand{\arraystretch}{1.17}
\newcommand{\code}[1]{\texttt{#1}}
\newcolumntype{Y}{>{\raggedright\arraybackslash}X}
\newcolumntype{L}[1]{>{\raggedright\arraybackslash}p{#1}}
\begin{document}
\begin{center}
{\Large\bfseries Mundo dos Blocos de Tamanho Variável}\\[5pt]
Fundamentos de Inteligência Artificial -- UFAM / ICOMP\\[6pt]
Camila dos Santos Soares; Daniel Nunes Santos; Nicoly Lima de Souza;\\
Pedro Matheus Melo Pereira; Pedro Vinícius Diaz de Alencar
\end{center}

\section{Introdução ao Problema}
O objetivo é mover blocos de comprimentos diferentes sem colisões e com apoio suficiente. Resolvemos a Situação 1 até Sf4, a Situação 2 até S5 e a Situação 3 até S7. Primeiro definimos o domínio e os planos manuais; depois usamos Python e MiniSat 2.2 para gerar planos a partir de cláusulas CNF.

\section{Descrição Formal do Mundo dos Blocos de Tamanho Variado}
\textbf{Domínio em LPO.} Os blocos são $B=\{a,b,c,d\}$, com comprimentos $\ell(a)=1$, $\ell(b)=1$, $\ell(c)=2$ e $\ell(d)=3$. A mesa é $T$. Existem sete pontos, de 0 a 6, e seis slots $s_i=[i,i+1]$, para $i=0,\ldots,5$. Um bloco em $p$ cobre os slots de $p$ até $p+\ell(b)-1$. A posição deve satisfazer $0\leq p\leq6-\ell(b)$; os níveis vão de 0 a 3, sendo 0 a mesa. Encostar pelas bordas não significa compartilhar um slot.

\begin{tabularx}{\linewidth}{lY}
\toprule Predicado & Significado no instante $t$ \\ \midrule
$at(b,p,t)$ & O bloco $b$ começa no ponto $p$.\\
$lev(b,l,t)$ & O bloco $b$ está no nível $l$.\\
$clr(b,t)$ & Nenhum bloco mais alto cobre slots de $b$.\\
$on(b,y,t)$ & $b$ tem apoio em $y$; pode haver mais de um apoio.\\ \bottomrule
\end{tabularx}

Cada bloco possui exatamente uma posição e um nível. Usamos a relação auxiliar de cobertura:
\[cov(b,s,l,t)\leftrightarrow lev(b,l,t)\land\exists p\,[at(b,p,t)\land p\leq s<p+\ell(b)].\]
Para blocos, $on(b,y,t)$ vale quando estão em níveis consecutivos e compartilham pelo menos um slot; $on(b,T,t)\leftrightarrow lev(b,0,t)$. Já $clr(b,t)$ vale quando não existe outro bloco em nível maior cobrindo algum slot de $b$. Essas condições se aplicam a todos os blocos e instantes do domínio.

\textbf{Estabilidade.} Para todo bloco acima da mesa, pelo menos $\lceil\ell(b)/2\rceil$ slots sob ele precisam estar ocupados no nível imediatamente inferior. Assim, $a$ e $b$ precisam de 1 slot, $c$ de 1 e $d$ de 2. Conta-se o conjunto de slots apoiados: $d$ pode formar uma ponte sobre $a$ e $b$ com um vão no meio.

\textbf{Ação.} $move(b,y,p,t)$ coloca $b$ sobre $y$ a partir de $p$. Exige $b$ livre, posição válida, destino diferente do atual, espaço no destino e acima dele, além de estabilidade após o movimento. Se $y$ é bloco, deve haver sobreposição horizontal e o novo nível é $nivel(y)+1$; para a mesa, é 0. Uma ação ocorre por passo, sem transportar outros blocos.

Para reproduzir os blocos lado a lado nas figuras, substituímos $clr(y)$ do manual por espaço livre na região utilizada do apoio. O bloco movido continua exigindo $clr$ e não pode entrar sob outro bloco, inclusive no vão de uma ponte.

\newpage
\begin{tabularx}{\linewidth}{lY}
\toprule Efeito & O que acontece \\ \midrule
Adds & Nova posição, novo nível e novos apoios; antigos apoios ficam livres se nada permanecer acima.\\
Deletes & Posição e nível antigos quando diferentes, apoios perdidos e $clr$ dos apoios agora cobertos.\\
Permanência & Posição e nível dos demais blocos; comprimentos, mesa, pontos e slots não mudam. $clr$ e $on$ acompanham a geometria.\\ \bottomrule
\end{tabularx}

\textbf{Ordem parcial.} \code{A -- P --> B} significa que A estabelece uma condição P necessária a B, preservada entre as ações. Chamando os quatro movimentos manuais da Situação 1 de A1 a A4, temos:
\begin{lstlisting}
A1 -- clr(a) --> A2
A2 -- slot s3 livre na mesa --> A3
A3 -- clr(c) --> A4
\end{lstlisting}
Movimentos sem dependência podem trocar de ordem. As pré-condições e a permanência representam essas dependências; o SAT produz uma sequência compatível, sem receber o plano manual como restrição.

Por exemplo, na Situação 2 automática, após colocar $c$ sobre $d$, os movimentos finais de $a$ e $b$ podem trocar de ordem, pois usam slots distintos de $c$.

\section{Codificação CNF}
Substituímos as variáveis da LPO pelos blocos, posições, níveis e instantes possíveis. Cada átomo vira uma variável booleana numerada. Usamos $at$, $lev$, $clr$, $move$ e duas auxiliares: $pos(b,p,l,t)$, que reúne posição e nível, e $occ(s,l,t)$, que indica slot ocupado. $on$ é reconstruído, sem variável própria. Os horizontes são $H=4,5,6$: estados de 0 a H e ações de 0 a H$-1$.

\begin{tabularx}{\linewidth}{L{3.5cm}Y}
\toprule Grupo de cláusulas & Regra codificada \\ \midrule
Inicial e meta & Fixam $at$ e $lev$ nos extremos do plano.\\
Unicidade e colisões & Uma posição e um nível por bloco; no máximo um ocupante por célula.\\
Estabilidade e clear & Apoio mínimo abaixo; $clr$ equivale à ausência de ocupação acima.\\
Pré-condições e efeitos & Bloco livre, apoio compatível, destino livre, nova posição e novo nível; sem ação vazia.\\
Frame axioms & Sem movimento de $b$, seus $at$ e $lev$ permanecem. A definição de $clr$ preserva seu valor enquanto a ocupação acima não mudar.\\
Ação única & Exatamente um $move$ por instante.\\ \bottomrule
\end{tabularx}

Exemplos de conversão: $move\rightarrow clr$ vira $\neg move\lor clr$. A equivalência $pos\leftrightarrow(at\land lev)$ é representada por três cláusulas:
\[(\neg pos\lor at)\land(\neg pos\lor lev)\land(\neg at\lor\neg lev\lor pos).\]
Para $d$, os três slots de apoio $x,y,z$ exigem $(x\lor y)\land(x\lor z)\land(y\lor z)$, condicionado à posição de $d$. Isso garante pelo menos dois slots ocupados. No DIMACS, os IDs representam variáveis, o sinal negativo indica negação e cada cláusula termina em \code{0}.

\newpage
\section{Exemplos dos 3 Cenários Codificados para CNF em LP}
Cada par $(p,l)$ representa $at(bloco,p,t)\land lev(bloco,l,t)$. Todos os pares de uma linha valem juntos: inicialmente em $t=0$ e na meta em $t=H$. As linhas fixam os átomos correspondentes com cláusulas unitárias.

\begin{center}
\begin{tabular}{lcccc}
\toprule Situação / estado & $a$ & $b$ & $c$ & $d$ \\ \midrule
1 e 3: S0 & (3,0)&(5,0)&(0,0)&(3,1)\\
1: Sf4 & (0,1)&(5,0)&(0,0)&(2,0)\\
2: S0 & (0,1)&(1,1)&(0,0)&(3,0)\\
2: S5 & (4,2)&(5,2)&(4,1)&(3,0)\\
3: S7 & (0,1)&(1,1)&(0,0)&(3,0)\\ \bottomrule
\end{tabular}
\end{center}

\textbf{Planos manuais.} Na tabela, $move(b,y,p)$ omite apenas o instante, igual ao passo menos 1. O par após a seta é a nova posição e nível do bloco movido; os demais permanecem iguais. Assim, cada linha permite obter o próximo estado até a meta.

\begin{center}\small
\begin{tabular}{clll}
\toprule Passo & Situação 1 & Situação 2 & Situação 3\\ \midrule
1 & $move(d,c,0)\to(0,1)$ & $move(b,T,2)\to(2,0)$ & $move(d,c,0)\to(0,1)$\\
2 & $move(a,b,5)\to(5,1)$ & $move(a,b,2)\to(2,1)$ & $move(a,b,5)\to(5,1)$\\
3 & $move(d,T,2)\to(2,0)$ & $move(c,d,4)\to(4,1)$ & $move(d,T,2)\to(2,0)$\\
4 & $move(a,c,0)\to(0,1)$ & $move(a,c,4)\to(4,2)$ & $move(a,c,0)\to(0,1)$\\
5 & -- & $move(b,c,5)\to(5,2)$ & $move(b,c,1)\to(1,1)$\\
6 & -- & -- & $move(d,T,3)\to(3,0)$\\ \bottomrule
\end{tabular}
\end{center}

Nas Situações 1 e 3, $a$ vai temporariamente sobre $b$: deixá-lo em $p=4$ na mesa, como no exemplo do manual, impediria $d$ de ocupar $[2,5]$. Na Situação 3, S5 e S6 desenham a mesma configuração; não é necessário um movimento vazio.

\textbf{Comparação com as execuções reais.}

\begin{tabularx}{\linewidth}{lL{2.2cm}Y}
\toprule Situação & Manual / solver & Comparação \\ \midrule
1 & 4 / 4 ações & Mesma sequência, chegando a Sf4.\\
2 & 5 / 5 ações & Outro apoio temporário para $a$ e inversão dos dois últimos movimentos; mesmo S5.\\
3 & 6 / 6 ações & Mesma sequência, chegando a S7.\\ \bottomrule
\end{tabularx}

Na Situação 2, a sequência do solver foi: $move(b,T,2)$; $move(a,d,3)$; $move(c,d,4)$; $move(b,c,5)$; $move(a,c,4)$. Os três resultados são SAT. Os horizontes vêm dos planos manuais; não se afirma optimalidade.

No início das Situações 1 e 3, $on(d,a,0)$ e $on(d,b,0)$ descrevem a ponte. Em Sf4, $a$ fica sobre $c$; em S5, $a$ e $b$ ficam sobre $c$, e $c$ sobre $d$; em S7, $a$ e $b$ ficam sobre $c$. Os blocos restantes ficam na mesa.

\newpage
\section{Mapeamento: Descrição Formal para Código}
\begin{tabularx}{\linewidth}{L{5cm}Y}
\toprule Conceito & Parte do código \\ \midrule
Blocos, comprimentos e cenários & \code{BLOCKS}, \code{MAX\_POINT}, \code{MAX\_LEVEL}, \code{CENARIOS} em \code{bw2cnf\_var.py}\\
Posições, slots e sobreposição & \code{posicoes}, \code{slots}, \code{sobrepoe}\\
Predicados, ações e regras CNF & \code{gerar\_cnf}; \code{var} numera os símbolos\\
DIMACS e MAP & \code{salvar}\\
MiniSat 2.2 e resposta numérica & \code{resolver}\\
Plano legível e apoios & \code{interpretar} e \code{derivar\_on}, em \code{interpretar.py}\\ \bottomrule
\end{tabularx}

\section{Execução Passo a Passo: Gerar o CNF, Mapeamento e Execução}
No Linux, abra o terminal dentro de \code{TrabalhoIA\_MundoBloclos}. Prepare o ambiente uma vez:
\begin{lstlisting}
python3 -m venv ../.venv_blocos
source ../.venv_blocos/bin/activate
python3 -m pip install -r requirements.txt
\end{lstlisting}
Em Ubuntu/Mint, se faltar venv, instale \code{python3-venv} com o gerenciador de pacotes. O ambiente fica fora da pasta do trabalho. Gere CNF e MAP, execute MiniSat e interprete os três resultados:
\begin{lstlisting}
for i in 1 2 3; do
    python3 bw2cnf_var.py "$i" --resolver
    python3 -B interpretar.py "situacao$i/resultado$i.txt"
done
\end{lstlisting}
O argumento 1, 2 ou 3 seleciona a situação. Sem \code{--resolver}, o primeiro comando apenas gera CNF e MAP. Com essa opção, lê o CNF salvo e executa MiniSat 2.2 pela biblioteca \href{https://pysathq.github.io/docs/html/api/solvers.html}{PySAT}, sem instalar outro executável. Cada pasta \code{situacaoN} recebe seu CNF, MAP e \code{resultadoN.txt}. A opção \code{-B} evita cache do Python. Em um novo terminal, reative o ambiente com \code{source ../.venv\_blocos/bin/activate}.

\section{Interpretação da Saída do SAT Solver}
\code{resultadoN.txt} contém \code{SAT} e o modelo real do solver: números positivos são variáveis verdadeiras e negativos são falsas. \code{UNSAT} indicaria ausência de plano no horizonte usado. O MAP associa cada número a um símbolo, por exemplo, \code{move(d,c,0,0)}.

\code{interpretar.py} lê o MAP da mesma pasta, seleciona as ações verdadeiras e ordena por $t$. Exemplo: \code{t=0: mover 'd' para cima de 'c', no ponto p=0.} Depois mostra posições e níveis finais e reconstrói $on$: nível 0 indica mesa; nos demais níveis, procura todos os blocos no nível imediatamente inferior com slots compartilhados. Assim, reconhece um apoio simples ou uma ponte. Não se deve misturar o resultado de uma situação com o MAP de outra.
\end{document}
